\section{Methodology}

\subsection{Bubble-overlap network construction}

At each date $t$, we construct a directed graph where nodes represent firms and edges represent temporally ordered bubble overlap.
An edge $i \rightarrow j$ is added when firm $i$ is in an active bubble and firm $j$ enters a bubble later within that overlap period.
Node features include normalized $\Delta$CoVaR and bubble duration (days), with optional eigenvector centrality.

\subsection{Temporal graphs and centrality analytics}

The pipeline constructs a time series of PyTorch Geometric graph snapshots (\texttt{results/temporal\_graphs.pkl}).
Centrality measures (degree, betweenness, eigenvector) are computed per snapshot and visualized as:
\begin{itemize}
  \item \texttt{documents/figures/fig\_centrality\_dynamics.png}
  \item \texttt{documents/figures/fig\_centrality\_heatmap.png}
\end{itemize}

\subsection{FRM dynamic networks (optional; heavier compute)}

We optionally construct rolling-window FRM networks via quantile regression to estimate tail dependence.
This step can be computationally expensive and is exposed via a dedicated flag in the build script.

\subsection{TGNN forecasting (optional)}

We provide a TGNN forecasting pipeline (GConvGRU/A3TGCN/EvolveGCN variants) that predicts next-step eigenvector centrality targets.
This module is optional and can be enabled for full builds.

